\chapter*{Preface}
\addcontentsline{toc}{chapter}{Preface}

These notes are a rotation foundation for understanding clustered spiking networks through mesoscopic theory. The central problem is to connect three descriptions of the same circuit: microscopic spike trains, cluster-level population variables, and low-dimensional stochastic dynamics.

The main path begins with the Litwin-Kumar and Doiron clustered balanced network \parencite{litwinKumarDoiron2012}, then builds the population-theory machinery needed to read finite-size mesoscopic equations \parencite{schwalgerDegerGerstner2017}. The later chapters use the Rostami et al.\ E/I clustered attractor model \parencite{rostami2024spiking} as the target architecture for a possible mesoscopic reduction. The SSN and variability-quenching chapters are included as a side branch because they offer a different explanation of stimulus-dependent variability \parencite{rubinVanHooserMiller2015,hennequinEtAl2018,churchlandEtAl2010}.

The notes deliberately avoid a polished-paper voice where uncertainty would be misleading. When a paper-specific parameter, claim, or manuscript status has not been checked, the source text leaves a visible citation or verification TODO. Those TODOs are part of the research artifact: they mark exactly where the next reading pass should replace a mechanism-level statement by a sourced one.

The practical objective is not to memorize every equation. It is to learn what each variable represents, how finite size enters, which statistics identify metastability, and what would distinguish noise-driven switching among stable attractors from deterministic mesoscopic irregularity.

\bigskip
\noindent\textit{Berlin, May 2026}
